% ---------------------------------------------------------------------------
% Answers the proposal's question in the terms the proposal asked it.
% ---------------------------------------------------------------------------

\section{Conclusion}
\label{s:sec:conclusion}

The question this work was proposed to answer was whether observability lock-in is a
syntactic accident or a semantic necessity. It is both, and which one it is depends on
which way you are going.

Leaving a proprietary vendor for an open one is a syntactic accident. We built the
translator the proposal described --- two inferred grammars, both admitting deterministic
one-token-lookahead parsers, a syntax-directed lowering into an intermediate
representation, an emission into PromQL, soundness stated relative to an ingestion mapping
because the two systems do not store the same thing --- and it translates $96.6\%$ of the
$8{,}770$ Datadog metric query occurrences in our corpus. The column we had reserved for
constructs beyond the target language's reach emptied of everything we can demonstrate,
and most of what emptied it were corrections to our own claims. $14$ occurrences remain
undecided and are reported as such. On this evidence, forward lock-in is work nobody has
done rather than a property of the languages.

Coming back is a semantic necessity, and the necessity is provable. A Datadog tag filter
is a glob, and globs are a proper subclass of the regular languages in a way that bites:
a glob language is a singleton or it is unbounded in length, no glob denotes a finite
language with two or more words, none denotes $\Sigma^{+}$, none denotes $A^{+}$ for a
proper non-empty $A\subsetneq\Sigma$, and an infinite language that omits a character of
the alphabet is not a union of finitely many glob languages at all. Whether a regular
expression has a glob is decidable, by a short argument about shortest words, with a
witness when one exists; run over all $184$ distinct concrete patterns in the corpus, that
procedure agrees with the syntax-directed proposer on every one. Of the $84$ refused
occurrences, $74$ are refused by a theorem about single globs together with a branch bound
we chose and state, and $10$ by a theorem about globs outright. No implementation effort
removes those ten.

So the lock-in is real, and it is not where the question assumed it would be. It is not
that a proprietary language says things an open one cannot; on this corpus the open
language said everything we could demonstrate the proprietary one saying. It is that the
proprietary filter language is weaker, so the migration is one-way: the queries that
survive the trip out are not all of the queries that could make the trip back. Leaving
costs work; returning costs expressiveness.

Two directions follow from this and are not pursued here. Whether glob-union-expressibility
is decidable without a bound on the number of branches is open, and the shortest-word
argument does not transfer to it. And the complete paper's finding that a dashboard query
does not determine its own meaning suggests the measurement that matters next is not of
query languages at all but of the resolution layer above them, which belongs to neither
vendor and which a common query language would not reach.
